% sections/1_intro.tex
\section{Introducción}

La resolución numérica de ecuaciones diferenciales en derivadas parciales es fundamental para el estudio de sistemas físicos complejos. En este contexto, la ecuación de Poisson gobierna fenómenos electrostáticos, donde el potencial eléctrico se relaciona con la distribución de carga. La determinación del potencial ante geometrías arbitrarias o condiciones de contorno específicas requiere el empleo de métodos numéricos cuando no existe solución analítica.

En el presente informe, se aborda la resolución computacional de la ecuación de Poisson en diversos escenarios. En primer lugar, se estudia la simetría cilíndrica de un sistema tipo condensador con tapas, resolviendo la ecuación en dos dimensiones ($r, z$) debido a la independencia angular del problema. Posteriormente, se presenta un sistema de trabajo automatizado para geometrías arbitrarias. Este enfoque procesa imágenes de entrada mediante la librería OpenCV (\texttt{cv2}) para detectar límites que pueden interpretarse como líneas de potencial constante o regiones de densidad de carga. A partir de estas condiciones topológicas, se resuelve la ecuación de Poisson numérica para diferentes tolerancias, ilustrando el método con ejemplos como el logotipo de la Facultad de Ciencias de la UAM y una foto de amigos en Hamburgo. Finalmente, se contrasta este método con el enfoque estándar para el ejercicio académico 9.1 del libro \citet{newman2012computational}, demostrando cómo el método basado en visión artificial reproduce fielmente los resultados esperados al interpretar un dibujo conceptual como densidades de carga.
